\section{Introduction}\label{sec:intro}

Modern Linux filesystems --- ext4, btrfs, bcachefs, and out-of-tree
ZFS --- have converged on \emph{detection} of metadata corruption
through cryptographic or polynomial checksums. \emph{Correction}
of detected corruption, when offered at all, is delegated either to
storage redundancy (RAID, mirrors) or to manual recovery tooling.
For radiation-prone deployments --- aerospace, certain industrial
control systems, high-altitude long-duration platforms --- this
operational model is insufficient. Single-event upsets (SEU) and
correlated burst errors require \emph{in-place} forward error
correction at the filesystem layer, layered with strict
fail-closed semantics that refuse to mount or proceed when the
trust chain breaks.

The FTRFS concept was introduced in 2015 by Fuchs, Langer and
Trinitis~\cite{fuchs2015ftrfs}. The work presented here is a
separate implementation, eleven years later, against a substantially
different technical context: the Linux kernel now ships a mature,
production-grade Reed--Solomon library
(\texttt{lib/reed\_solomon}~\cite{kernelrslib}) used by MTD, DVB
and EDAC subsystems; reproducible builds via Yocto LTS releases
are the norm in regulated embedded deployment; MIL-STD-882E
hazard tracking practice~\cite{milstd882e} has been formalized
for software-intensive systems; and post-quantum integrity
primitives (FIPS~205~\cite{fips205}) are entering the
standardization pipeline. FTRFS as documented here is a 2026
response to the 2015 problem statement, not a port.

\paragraph{Contributions.} This report documents the following
contributions, each verifiable against the public source tree at
\url{https://github.com/roastercode/FTRFS}:
\begin{itemize}
\item Universal Reed--Solomon FEC on inode (item~1, commit
  \texttt{ee6b6ae}) and superblock (item~2, commits
  \texttt{4227354}/\texttt{874dfdd}/\texttt{584440e}) metadata.
\item Strict layered defense: CRC32 detection followed by RS
  correction on CRC failure, with re-verification after correction.
\item Persistent radiation event journal: a 64-entry ring buffer
  in the superblock, with each entry CRC32-protected.
\item Fail-closed mount: strict version, scheme and incompatible
  feature checks; unknown structural bits cause mount refusal.
\item Single-source-of-truth on-disk layout: \texttt{offsetof()}
  and \texttt{BUILD\_BUG\_ON} sentinels enforce kernel/user-space
  symmetry at compile time (item~4a, commit \texttt{c75a067}).
\end{itemize}

\paragraph{Scope.} The format v4 bump (item~4) is designed and
frozen but not yet implemented; v2 of this report will document it.
Stage~4 (universal data-block protection) and stage~5 (offensive
security review) are described as future work.
